
%In this section, we provide the necessary background on grammar-guided LLMs generation.

% \noindent{\bf Notation.} 
Let the alphabet $\Sigma$ be a finite set of characters and $\epsilon$ denotes an empty string.
Given a set $S$, $S^i$ denotes the set of all $i$-length sequences that can be formed by concatenating elements from $S$,
and $S^* = \bigcup_{i \in \mathbb{N}} S^i$.
$\alphabets^{*}$ represents the set of all strings over characters in $\alphabets$, including the empty string $\epsilon$. 
% Further, we use $\alphabets^{+}$ to denote $(\alphabets^{*} - \epsilon)$.
% Given two strings $w_1, w_2 \in \alphabets^*$, we use $w_1.w_2$ to denote string obtained by concatenating $w_2$ to $w_1$.

% \begin{wrapfigure}{R}{0.53\textwidth}
\vspace{-.1in}
\begin{minipage}{0.55\textwidth}
\begin{algorithm}[H]
\small
\caption{Constrained LLM Generation}
\label{alg:masked}
%
\begin{flushleft}
\textbf{Inputs:} $M$: LLM, $\tokenizer$: tokenizer, $O_0$: input prompt string, \\
$f_m$: function that generates mask, $n_\textit{max}$: maximum generated tokens, $D$: any decoding algorithm \\
\textbf{Output:} string $O_n$
\end{flushleft}

\begin{algorithmic}[1]
\Function{ConstrainedGenerate}{$M$, $\tokenizer$, $f_m$, $O_0$}
\State $\curtokens \gets \text{Tokenize}(\tokenizer, O_0)$
\For{$i \in \{1, \dots n_\textit{max} \}$}
\State $\textit{scores} \gets M(\curtokens)$
\State $m \gets f_{m}(\curtokens, \tokenizer)$
\State $\textit{scores} \gets m \odot \textit{scores}$
\State $t_i \gets D(\textit{scores})$
\If{$t_i =$ EOS}
\State break
\EndIf
\State $\curtokens \gets \text{append}(\curtokens, t_i)$
\EndFor
\State $O_n \gets \text{Detokenize}(\tokenizer, \curtokens)$
\State \Return $O_n$
\EndFunction
%\item[]
\end{algorithmic}
\end{algorithm}

\end{minipage}
\vspace{-0.1in}
\end{wrapfigure}


\subsection{Language Models}

Current autoregressive language models (LM) operate on vocabulary $V \subseteq \Sigma^*$ of tokens.
A tokenizer takes an input prompt $O_0 \in \Sigma^*$,
which is a sequence of characters,
as input and converts $O_0$ into a sequence of tokens $t_1, t_2, \dots, t_k$. 
In order to generate the next token, the LM $M: V^* \to \mathbb{R}^{|V|}$ takes as input the sequence of tokens $t_1, t_2, \dots, t_k$, and outputs a vector of scores $\mathcal{S}$ over the vocabulary: $\mathcal{S} = M(t_1, t_2, \dots, t_k)$.
The softmax function $\textit{softmax}(\mathcal{S}_i) = \exp(\mathcal{S}_i)/\sum_j(\exp(\mathcal{S}_j))$ 
transforms $\mathcal{S}$ into a probability distribution over the vocabulary $V$, and then $t_{k+1}$ is sampled as the next token.

\sloppypar
\noindent{\bf Decoding.} 
% The language model $M$ is recurrently applied to generate a sequence of tokens $t_1, t_2 \dots t_k$. 
% When choosing the $(k + 1)$-th token, the probability distribution for the next token is obtained through $\text{softmax}(M(t_1, t_2 \dots t_k))$. 
Various approaches for token selection from this distribution have been explored in the literature such as greedy decoding, sampling, and beam search.
Each technique is repeated until the prediction of a special end-of-sequence token, EOS, or another stopping criterion is fulfilled. 
This iterative process is equivalent to sampling from a distribution over $V^*$, potentially resulting in multiple feasible decoding outputs.

\noindent{\bf Constrained Masking.} 
In the context of decoding, we encounter scenarios where excluding specific tokens at particular positions becomes crucial (e.g., excluding harmful words). 
This implies we can disregard these tokens and proceed with decoding based on the remaining set. 
An algorithm for such masking relies on a function $f_m$ to generate the mask $m$ based on the exact use case. 
In the mask $m \in \{0, 1\}^{|V|}$, $1$ indicates a viable token, and $0$ signifies a discarded one. 
Decoding methods mentioned earlier can be applied to $m \odot \textit{softmax}(\mathcal{S})$, where $\odot$ represents element-wise multiplication. 
% The resultant vector should be scaled by $1/\sum_i(m \times \textit{softmax}(\mathcal{S}))_i$ to restore correct probabilities.
% Algorithm~\ref{alg:masked} presents the steps for constrained decoding.

\subsection{Grammar-guided generation}

\noindent\textbf{Grammar:}
A formal language's syntax is defined by grammar, which comprises a set of production rules that specify all possible strings within that language. 
A grammar includes terminal and nonterminal symbols. Terminal symbols represent the actual characters or tokens; nonterminal symbols serve as placeholders that define patterns or structures within the language.
Most programming languages can be described using context-free grammar, which consists of production rules that apply to nonterminal symbols independently of their context. 
Each production rule is of the form $S \to S_1, S_2 \dots S_n$ with $S$ a single nonterminal symbol, and $S_1, S_2 \dots S_n$ a string of terminals and nonterminals. 
Single nonterminal $S$ on the left-hand side can be replaced by $S_1, S_2 \dots S_n$ \mbox{on the right-hand side.}


\noindent\textbf{Shift-Reduce LR Parser:}
An LR parser is a bottom-up parser used for analyzing context-free grammars (CFGs)~\cite{aho86}. 
It handles deterministic grammars by reading input from left to right, constructing a rightmost derivation in reverse (hence LR). 
The parser uses a shift-reduce method, shifting symbols onto a stack until a sequence matches a grammar rule. 
When a match is found, the symbols on the stack are reduced by applying the rule, replacing them with the corresponding non-terminal. 
This process repeats until the entire input is successfully parsed or an error occurs.

\noindent\textbf{Constrained grammar-guided generation:}
 Recent works have explored constrained grammar-guided LLM generation~\cite{Wei_2023, 10.1145/3591300, guidance, willard2023efficient, scholak-etal-2021-picard, poesia2022synchromesh, geng2023grammar, beurerkellner2024guiding, ugare2024syncodellmgenerationgrammar}. 
These methods typically incorporate an incremental parser alongside the LLM, which parses the partial output at each decoding step. 
The parsing results are then used to filter out tokens that would lead to syntactically invalid sequences.
